\subsection{Conclusion}

Minimax regret identifies the equipment choice whose largest shortfall from the scenario benchmark is smallest. We use the comfort ranking as a benchmark. The Grove Hall example shows how local limits on purchase cost, electricity, and water alter the options available to minimax.

The scenario benchmark matters because the hour with the highest discomfort can differ from the hour with the largest regret. Discomfort measures how far SET remains above the comfort target. Regret measures how much of that discomfort another feasible device could have avoided. A severe hour can leave substantial discomfort under every device and still produce a small gap between them. An hour with a lower baseline can determine the equipment choice if one option delivers substantially more relief than the others. Tracking both quantities shows where equipment choice has the greatest effect and where the venue faces a larger cooling problem.

We also distinguish a purchase decision from an operating decision. A purchase commits the venue to equipment that must be useful across the selected weather cases. During operation, staff can respond to the conditions that occur in a given hour. They can leave equipment off when the baseline meets the target and use the resource estimates to weigh additional cooling when discomfort remains. The next section reports the purchase comparison across all twelve venues.
